\section{Introduction and Related Work Draft}
\label{sec:rw-draft}

\paragraph{Problem and motivation.}
We consider first-order methods for smooth nonconvex optimization with
nonlinear equality constraints.  Augmented Lagrangian methods provide a
natural mechanism for coupling stationarity and feasibility, but their
nonconvex complexity analysis is delicate: large penalty parameters improve
feasibility while worsening conditioning, and proximal augmented Lagrangian
schemes often obtain their guarantees through inner subproblems that must be
solved with increasing accuracy.  This tension motivates single-loop
linearized schemes whose iterations remain explicit, while still retaining a
provable mechanism for controlling both the multiplier sequence and the
constraint residual.

\paragraph{Our perspective.}
The proposed DDLALM framework introduces a dual proximal damping term and
eliminates the damped multiplier in closed form.  This produces a reduced
augmented Lagrangian with an effective penalty
$\bar\rho_k=\rho+s_k$, while the anchor multiplier $\bmu^k$ evolves through a
damped dual update.  The resulting method separates two roles that are often
entangled in standard ALM analyses: the effective penalty controls feasibility,
whereas the anchor multiplier is the stable dual object used in the descent
argument.  Under a compact working-set assumption and a uniform LICQ condition,
we prove a weighted finite-length estimate; with
$\bar\rho_k\asymp k^{1/3}$ and $\eta_k\asymp k^{-1/3}$, this yields the
standard $O(\epsilon^{-3})$ complexity for finding an approximate KKT point.

\paragraph{Penalty and proximal augmented Lagrangian methods.}
Worst-case complexity for nonconvex constrained optimization has been studied
extensively through penalty, proximal-point, and augmented-Lagrangian
frameworks.  For nonlinear equality constraints, the proximal augmented
Lagrangian method of \citet{xie2021proximal} is a central reference point:
the proximal term stabilizes the subproblem sequence and enables
non-asymptotic complexity guarantees.  Related inexact augmented Lagrangian
schemes, including \citet{li2021rate} and
\citet{sahinInexactAugmentedLagrangian2019}, further clarify how approximate
subproblem solutions and penalty growth affect stationarity and feasibility.
Quadratic penalty methods offer another classical route, and recent analyses
with adaptive subproblem accuracy show that feasibility-aware stopping rules
can reduce unnecessary inner work when the iterate is far from feasible.
These methods provide strong theoretical baselines, but they still rely on a
penalty or augmented-Lagrangian subproblem viewpoint.  In contrast, our method
keeps each primal step fully linearized and explicit, and the effective
penalty appears through dual damping rather than through a conventional
outer-loop penalty update.

\paragraph{Linearized ALM and single-loop primal-dual methods.}
Single-loop first-order methods avoid expensive nonconvex subproblem solves,
but their analysis must compensate for the loss of exact subproblem
minimization.  \citet{kimNewLagrangianbasedFirstorder2023} introduced a
Lagrangian-based first-order method for nonlinear equality constraints using a
proximal-perturbed dual structure that yields bounded primal-dual iterates
without directly assuming bounded multipliers.  In stochastic settings,
\citet{alacaoglu2024complexity} studied single-loop algorithms for nonlinear
programming with stochastic objective and constraints, while
\citet{shi2025momentum} developed a momentum-based linearized augmented
Lagrangian method for stochastic nonconvex constrained optimization.  These
works show that linearization can be made theoretically meaningful, but the
main complexity objectives are different from ours: several results address
stochastic oracle complexity, expectation constraints, or broader functional
constraints.  Our analysis focuses on the deterministic nonlinear equality
case and derives the complexity from a reduced-potential finite-length
mechanism induced by dual damping.

\paragraph{Proximal-perturbed Lagrangian frameworks.}
The proximal-perturbed Lagrangian viewpoint has recently emerged as a useful
way to modify the geometry of the dual variables.  Kim's false-penalty and
P-Lagrangian constructions \citep{kimLagrangianBasedMethodFalse2023,
kimNewLagrangianbasedFirstorder2023} introduce strong concavity in the dual
variables and help avoid explicit bounded-multiplier assumptions.  More
recently, \citet{anonymous2026beyond} proposed proximal-perturbed augmented
Lagrangian methods for nonconvex functional constrained optimization and
obtained a single-loop first-order method with
$\widetilde O(\epsilon^{-2})$ iteration complexity under their functional
constraint setting.  Our method is aligned with this general principle, but
uses the perturbation differently: the dual damping variable is eliminated
to expose an effective penalty schedule $\bar\rho_k=\rho+s_k$, and the proof
tracks a stable anchor multiplier separately from the virtual multiplier.
This distinction is central to the descent and residual analysis.

\paragraph{Stochastic constrained optimization.}
Stochastic constrained optimization provides another perspective on why
penalty growth is costly.  In semi-stochastic settings, variance-reduction and
momentum mechanisms have been used to reach improved oracle complexities;
representative examples include stochastic inexact ALM
\citep{li2024stochastic}, momentum-based linearized ALM
\citep{shi2025momentum}, and variance-reduced first-order methods for
deterministically constrained stochastic nonconvex optimization
\citep{lu2026variance}.  Adaptive directional decomposition methods further
organize the search direction into objective-reduction and
constraint-reduction components \citep{shi2025adaptive}.  These papers pursue
stochastic oracle guarantees, whereas the present work is deterministic.
Nevertheless, they highlight the same structural difficulty: if feasibility
is enforced mainly through a large merit or penalty parameter, then the
smoothness constants and variance terms can become the dominant bottleneck.
DDLALM isolates this issue in the deterministic setting by using a scheduled
effective penalty together with a damped multiplier anchor.

\paragraph{ADMM and Bregman-type variants.}
ADMM and proximal ADMM variants are closely related to augmented Lagrangian
methods, especially for linearly constrained and block-separable nonconvex
composite problems.  Recent work on dual-descent ALM/ADMM
\citep{sun2024dual_descent_alm_admm}, proximal-perturbed ADMM, and Bregman
ADMM studies how proximal or Bregman regularization stabilizes nonconvex
block updates.  These methods are most natural when the model has linear
constraints or a separable block structure.  Our target problem has nonlinear
equality constraints and does not rely on an ADMM splitting.  We therefore use
this literature as methodological background, while positioning DDLALM as a
linearized augmented-Lagrangian method for nonlinear equality-constrained
nonconvex optimization.

\paragraph{Positioning.}
The closest conceptual neighbors are proximal ALM for nonlinear equality
constraints, P-Lagrangian methods that regularize the dual geometry, and
linearized ALM methods that avoid inner minimization.  DDLALM combines these
threads in a different way: it keeps the primal step explicit, introduces a
virtual multiplier through dual proximal damping, updates only a damped
anchor multiplier, and analyzes the method through a reduced potential with a
scheduled effective penalty.  This gives a clean deterministic route to an
$O(\epsilon^{-3})$ approximate-KKT complexity while preserving a genuinely
single-loop first-order iteration.
